%-------------------------------------------------------------------------------
%	SECTION TITLE
%-------------------------------------------------------------------------------
\cvsection{Formazione}

%-------------------------------------------------------------------------------
%	CONTENT
%-------------------------------------------------------------------------------
\begin{cventries}

%---------------------------------------------------------
  \cventry
    {Laurea Magistrale in Ingegneria dell'Intelligenza Artificiale e della Sicurezza Informatica{\enskip\cdotp\enskip}110 e Lode}
    {Università degli Studi di Enna "Kore"}
    {Enna, Italia}
    {2022 - 2024}
    {
      \begin{cvitems}
        \item {Specializzazione in \textbf{Machine Learning e Artificial Intelligence}: progettazione e implementazione di modelli di \textbf{classification, regression e clustering} utilizzando \textbf{Python}.}
        \item {Esperienza con principali algoritmi di ML: \textbf{logistic regression, support vector machines, neural networks, gaussian mixture models, k-means clustering}.}
        \item {Applicazione di tecniche di \textbf{dimensionality reduction} (\textbf{PCA, LDA}) e valutazione dei modelli tramite metriche statistiche e decision theory.}
        \item {Progettazione di \textbf{sistemi intelligenti}: agenti intelligenti, algoritmi di ricerca, rappresentazione della conoscenza, \textbf{Bayesian networks} e sistemi multi-agente.}
        \item {Competenze in \textbf{Human-Centered AI}: progettazione di sistemi AI centrati sull’utente, valutazione della \textbf{user experience}, usabilità, sicurezza e aspetti etici dell’intelligenza artificiale.}
        \item {Esperienza in \textbf{Speech Processing e Computational Linguistics}: speech recognition, speech synthesis, speaker recognition, signal processing e modelli statistici per il linguaggio.}
        \item {Competenze avanzate in \textbf{Cybersecurity}: autenticazione, access control, sicurezza di reti e sistemi operativi, sicurezza dei database e protocolli di sicurezza.}
        \item {Analisi di \textbf{cyber attacks e vulnerabilità di rete}, gestione degli incidenti, tecniche di \textbf{ethical hacking} e protezione delle infrastrutture IT.}
        \item {Studio di \textbf{AI for Security}: adversarial attacks nei modelli di machine learning e utilizzo dell’AI per il rilevamento di minacce informatiche.}
        \item {Progettazione di \textbf{sistemi biometrici}: riconoscimento facciale, impronte digitali, riconoscimento dell’iride e sistemi di autenticazione multi-biometrica.}
      \end{cvitems}
    }

%---------------------------------------------------------
  \cventry
    {Laurea Triennale in Ingegneria Informatica{\enskip\cdotp\enskip}96 / 110}
    {Università degli Studi di Enna "Kore"}
    {Enna, Italia}
    {2019 - 2022}
    {
      \begin{cvitems}
        \item {Solida formazione in \textbf{Computer Science fundamentals}: algoritmi, strutture dati, architettura dei calcolatori e progettazione di algoritmi per la risoluzione di problemi ingegneristici.}
        \item {Sviluppo software in \textbf{C, Java e Python} con progettazione \textbf{object-oriented}, modellazione tramite \textbf{UML} e implementazione di applicazioni modulari e riutilizzabili.}
        \item {Progettazione e gestione di \textbf{database relazionali}: modellazione concettuale e logica, normalizzazione e interrogazione con \textbf{SQL/MySQL (DDL, DML)}.}
        \item {Competenze in \textbf{Operating Systems e Concurrency}: gestione di processi e thread, sincronizzazione (mutex, semafori), IPC e programmazione di sistema in \textbf{C su Unix/Linux}.}
        \item {Fondamenti di \textbf{Computer Networks}: protocolli Internet, livelli di rete e trasporto, indirizzamento IP e analisi delle performance di rete.}
        \item {Introduzione a \textbf{Machine Learning e Artificial Intelligence} con sviluppo di modelli e applicazioni software nell’ecosistema \textbf{Python}.}
        \item {Principi di \textbf{Cybersecurity e crittografia}: algoritmi di cifratura, autenticazione e progettazione di soluzioni software sicure.}
        \\
      \end{cvitems}
    }
\cventry
    {Diploma di Maturità Classica}
    {Liceo Classico Tommaso Fazello}
    {Sciacca, Italia}
    {}
    {
      \begin{cvitems}
        \item {Formazione umanistica con sviluppo di \textbf{analisi critica, capacità comunicative e problem solving}.}
        \\
      \end{cvitems}
    }

%---------------------------------------------------------
\end{cventries}